\section{Conclusions}
We have presented a simulation-based framework for quantifying the impact of optical ghosts on LSSTCam imaging, motivated by the need to assess compliance with the LSST system requirements. We combine the Yale Bright Star Catalog, empirically calibrated magnitude transformations, and optical ray-tracing simulations using \texttt{Batoid}, to estimate both the morphology and surface brightness of the dominant ghosts and measure the resulting ghost-impacted area on the focal plane. These simulations assumed a 2\% reflectance for each of the optical surfaces, which approximately match  the system engineering simulations.

Applying this framework to $\sim$ 3,100 visits from the intermittent DRP (w37), we find that the average fraction of the Rubin imaging that is significantly impacted by ghosts is 0.57\% when averaged across all bands. This value is below the 1\% threshold specified by the LSST system requirement. The impacted area exhibits strong band dependence, with the largest contributions occurring in the {\it u} band due to its lower sky background and higher filter reflectance, and minimal impact in the {\it r} and {\it i} bands. It should be noted that the LSST requirement is not satisfied in the {\it u} band alone.  The impacted area is also heavily dependent on the field being observed, as individual visits containing very bright stars have substantially larger impacted fractions.

We further explored the dependence of ghost-impacted area on stellar magnitude and off-axis position, finding that star brightness is the dominant driver of ghost impact, with only a weak dependence on field angle. Using these results, we constructed all-sky maps of expected ghost-impacted area based on the Yale Bright Star Catalog. Overall, our results demonstrate that optical ghosts are a non-negligible source of contamination for low-surface-brightness science in LSST, particularly in the presence of very bright stars. 

We use the CBP to generate optical ghosts on the focal plane. 
We find good agreement between the morphology of the observed CBP ghosts and of those generated by \texttt{Batoid}.
We use \texttt{Batoid} to generate morphological templates for each of the most prominent ghosts and fit the fluxes of the ghosts to derive the reflectances of the optical surfaces.
We find that the reflectances of the optics generally agree with estimates generated by the systems engineering simulations, thus validating the assumption that was made when assessing ghost impact.